\documentclass[11pt]{article}
\usepackage[a4paper,margin=1in]{geometry}
\usepackage[T1]{fontenc}
\usepackage[utf8]{inputenc}
\usepackage{lmodern}
\usepackage{amsmath,amssymb}
\usepackage{microtype}
\setlength{\parindent}{1.2em}
\setlength{\parskip}{0.25em}
\newcommand{\dd}{\,d}
\begin{document}

\begin{center}
{\Large G. Lejeune Dirichlet's Werke. II.}\par\vspace{1.2em}
{\Large XVII.}\par\vspace{0.5em}
{\LARGE \textbf{Sur un théorème relatif aux séries.}}\par\vspace{1em}
{\large Par M. G. Lejeune Dirichlet}\par\vspace{0.5em}
{\small Liouville, Journal de Mathématiques pures et appliquées, Série II, Tome I, 1856, p. 80--81.}
\end{center}

\begin{center}
{\large \textbf{Sur un théorème relatif aux séries.}}
\end{center}

Le théorème que je vais établir d'une manière nouvelle est le même qui m'a servi comme lemme dans plusieurs Mémoires précédents. La démonstration que j'en ai donnée dans le premier de ces Mémoires\footnote{Journal de Crelle, tome XIX. Bd. I, S. 414--417 dieser Ausgabe von G. Lejeune Dirichlet's Werken. K.} suppose une certaine condition qui, bien qu'elle se trouve remplie dans la plupart des applications qu'on peut faire du lemme, n'est pas indispensable, comme j'en ai déjà fait la remarque ailleurs, et comme on pourra au reste le voir dans ce qui va suivre.

Je commence par un cas particulier très simple et qui consiste en ce que, \(a\) et \(b\) désignant des constantes positives et \(\rho\) une variable également positive, on a:
\[
\lim_{\rho\to0}\rho\sum_{n=0}^{n=\infty}\frac{1}{(b+na)^{1+\rho}}=\frac{1}{a},
\]
la limite se rapportant au décroissement indéfini de \(\rho\). Pour démontrer ce cas particulier, considérons l'intégrale:
\[
\int_b^\infty\frac{dx}{x^{1+\rho}}=\frac{1}{\rho b^\rho}
\]
comme l'aire d'une courbe ayant \(x\) pour abscisse; cette courbe se rapprochant de plus en plus de l'axe des \(x\), si l'on mène des parallèles à cet axe par les extrémités des ordonnées qui répondent à \(x=b,b+a,\) etc., on aura des rectangles extérieurs et intérieurs à notre aire, et l'on conclura sur-le-champ que la somme dont il s'agit d'avoir la limite, se trouve comprise entre les deux quantités:
\[
\frac{1}{ab^\rho}
\quad\text{et}\quad
\frac{1}{ab^\rho}+\frac{\rho}{b^{1+\rho}},
\]
dont la limite commune est \(1/a\).

Pour passer maintenant au théorème général, soient:
\[
k_1,\quad k_2,\quad k_3,\quad \ldots,\quad k_n,\quad \ldots,
\]
des constantes positives rangées par ordre de grandeur croissante, de sorte que \(k_n\leq k_{n+1}\), et supposons que ces constantes soient telles, qu'en désignant par \(t\) une variable continue et positive, et par \(T\) le nombre des termes de notre suite qui ne surpassent pas \(t\), le rapport:
\[
\frac{T}{t}
\]
converge vers la limite finie \(\alpha\) lorsque \(t\) croît indéfiniment. Cela supposé, je dis que cette même quantité \(\alpha\) sera aussi la limite de la somme:
\[
\tag{1} \rho\sum_{n=1}^{n=\infty}\frac{1}{k_n^{1+\rho}}
\]
pour une valeur infiniment petite de \(\rho\).

D'après la condition supposée, on peut, \(\delta\) désignant un nombre aussi petit que l'on voudra, en choisir un autre \(\tau\) assez grand pour que, \(t\) surpassant \(\tau\), on ait constamment:
\[
\alpha-\delta<\frac{T}{t}<\alpha+\delta.
\]
Cela fait, soit \(N\) la valeur de \(T\) qui répond à \(t=\tau\), et ne considérons d'abord que les termes de notre série ayant un indice \(n>N\). Un pareil terme \(k_n\) peut présenter deux cas, selon que l'on aura:
\[
k_n<k_{n+1}\quad\text{ou}\quad k_n=k_{n+1}.
\]

Dans le premier cas, la valeur de \(T\) qui répond à \(t=k_n\) sera \(n\), et l'on aura:
\[
\alpha-\delta<\frac{n}{k_n}<\alpha+\delta.
\]
Dans le second cas, soit, parmi les termes qui suivent \(k_n\), \(k_{m'}\) le dernier qui soit encore égal à \(k_n\); soit de plus, parmi les termes précédents, \(k_m\) le premier dont la valeur soit inférieure à celle de \(k_n\). Si maintenant nous faisons croître \(t\) à partir de \(t=k_m\), \(T\) restera invariablement égal à \(m\), tant que \(t\) n'atteindra pas la valeur \(k_{m'}\), et notre rapport \(T/t\) approchant ainsi indéfiniment de \(m/k_{m'}\), en même temps qu'il reste toujours supérieur à \(\alpha-\delta\), nous aurons:
\[
\frac{m}{k_{m'}}>\alpha-\delta.
\]
Mais comme par le premier cas nous avons aussi:
\[
\frac{m'}{k_{m'}}<\alpha+\delta,
\]
et en outre:
\[
m<n<m',\qquad k_n=k_{m'},
\]
nous concluons, comme précédemment:
\[
\alpha-\delta<\frac{n}{k_n}<\alpha+\delta.
\]

Cette double inégalité étant mise sous la forme:
\[
(\alpha-\delta)^{1+\rho}\frac{\rho}{n^{1+\rho}}
<
\frac{\rho}{k_n^{1+\rho}}
<
(\alpha+\delta)^{1+\rho}\frac{\rho}{n^{1+\rho}},
\]
si nous sommons depuis \(n=N+1\) jusqu'à \(n=\infty\), il viendra:
\[
(\alpha-\delta)^{1+\rho}\rho\sum_{n=N+1}^{n=\infty}\frac{1}{n^{1+\rho}}
<
\rho\sum_{n=N+1}^{n=\infty}\frac{1}{k_n^{1+\rho}}
<
(\alpha+\delta)^{1+\rho}\rho\sum_{n=N+1}^{n=\infty}\frac{1}{n^{1+\rho}}.
\]
Or l'expression \(\rho\sum 1/n^{1+\rho}\) rentrant dans le cas particulier examiné plus haut, en supposant:
\[
b=N+1,\qquad a=1,
\]
et ayant par suite l'unité pour limite, on voit que la partie de notre somme (1), qui s'étend depuis \(n=N+1\) jusqu'à \(n=\infty\), finira par rester comprise entre \(\alpha-\varepsilon\) et \(\alpha+\varepsilon\), le nombre \(\varepsilon\) étant tant soit peu supérieur à \(\delta\) et par conséquent, comme ce dernier, d'une petitesse arbitraire, et comme en même temps la première partie qui n'a qu'un nombre fini \(N\) de termes, converge évidemment vers zéro, il s'ensuit que la limite de l'expression (1) est en effet \(\alpha\); C. Q. F. D.
\end{document}
